\documentclass[12pt,a4paper]{report}
\setcounter{secnumdepth}{5} % setting level of numbering (default for "report" is 3).
\usepackage{natbib}
\usepackage[utf8]{inputenc}
\usepackage[french]{babel}
\usepackage[T1]{fontenc}
\usepackage{amsmath}
\usepackage{amsfonts}
\usepackage{amssymb}
\usepackage{graphicx}
\usepackage{appendix}
\usepackage{enumitem}
\usepackage{float} 
\usepackage[nottoc, notlof, notlot]{tocbibind}
\usepackage{pdfpages}
\usepackage{gensymb}
\usepackage{lipsum}
\usepackage[colorinlistoftodos,textwidth=3.7cm]{todonotes}
\usepackage[top=1in, bottom=1in, left=1in, right=1in]{geometry}
\usepackage{tocbibind}
\author{Sébastien Hervieu}

\newcounter{mypagecount}% create a new counter
\setcounter{mypagecount}{0}% set it to something just in case
\newenvironment{interlude}{% create a new environment for the unnumbered section(s)
	\clearpage
	\setcounter{mypagecount}{\value{page}}% use the new counter we created to hold the page count at the start of the unnumbered section
	\thispagestyle{empty}% we want this page to be empty (adjust to use a modified page style)
	\pagestyle{empty}% use the same style for subsequent pages in the unnumbered section
}{%
	\clearpage
	\setcounter{page}{\value{mypagecount}}% restore the incremented value to the official tally of pages so the page numbering continues correctly
}


\begin{document}
	
\newcommand{\para}{\paragraph*{}}
	
\newcommand{\todoRediger}[1]{\todo[color=red,inline]{REDIGER #1}}
\newcommand{\todoARevoir}[1]{\todo[color=red,inline]{A REVOIR #1}}
\newcommand{\todoCompleter}[1]{\todo[inline]{COMPLETER #1}}
\newcommand{\todoAmeliorer}[1]{\todo{AMELIORER: #1}}
\newcommand{\todoMiseEnForme}[1]{\todo[color=green,inline]{MISE EN FORME: #1}}
\newcommand{\todoCheck}[1]{\todo[color=cyan,fancyline]{VERIFIER #1}}
\newcommand{\todoAjouterRef}[1]{\todo[color=pink]{AJOUTER REF #1}}
\newcommand{\todoTrouverRefExt}[1]{\todo[color=red,inline]{TROUVER REF EXTERNE: #1}}
\newcommand{\todoObjectif}[1]{\todo[color=yellow,inline,size=\Large]{OBJECTIF: #1}}
\newcommand{\todoPeutEtre}[1]{\todo[color=teal,inline]{PEUT ETRE #1 ?}}

\begin{titlepage}

\newcommand{\HRule}{\rule{\linewidth}{0.5mm}} % Defines a new command for the horizontal lines, change thickness here

\center % Center everything on the page
 
%----------------------------------------------------------------------------------------
%	HEADING SECTIONS
%----------------------------------------------------------------------------------------

\textsc{\LARGE Université de Rennes 1}\\[1cm] 
\textsc{\Large }\\[0.5cm] % Major heading such as course name
\textsc{\large Master 2 Calcul scientifique et modélisation}\\
\textsc{Rapport de Stage}\\
%----------------------------------------------------------------------------------------
%	TITLE SECTION
%----------------------------------------------------------------------------------------

\HRule \\[0.4cm]
{ \huge \bfseries Etude et développement d’outils mathématiques pour estimer, en temps réel, le tassage et le volume d’un silo de maïs à partir de capteurs embarqués}\\[0.4cm] 
\HRule \\[1.5cm]
 
%----------------------------------------------------------------------------------------
%	AUTHOR SECTION
%----------------------------------------------------------------------------------------

\begin{minipage}{0.4\textwidth}
\begin{flushleft} \large
\emph{Auteur:}\\
Sébastien \textsc{Hervieu}
\end{flushleft}
\end{minipage}
~
\begin{minipage}{0.4\textwidth}

\begin{flushright}
	\emph{Tuteur de Stage:} \\
	Geoffroy \textsc{Etaix}
\end{flushright}

\begin{flushright} \large
	\emph{Tuteur Universitaire:} \\
	Fabrice \textsc{Mahé} 
\end{flushright}

\end{minipage}\\[1cm]

% If you don't want a supervisor, uncomment the two lines below and remove the section above
%\Large \emph{Author:}\\
%John \textsc{Smith}\\[3cm] % Your name

%----------------------------------------------------------------------------------------
%	DATE SECTION
%----------------------------------------------------------------------------------------

{\large \today}\\[1cm] % Date, change the \today to a set date if you want to be precise

%----------------------------------------------------------------------------------------
%	LOGO SECTION
%----------------------------------------------------------------------------------------
\includegraphics[height=3cm]{img/logo-tellusenv.png} \\
\includegraphics[height=3cm]{img/univ.jpeg}\\[1cm] % Include a department/university logo - this will require the graphicx package
 
%----------------------------------------------------------------------------------------

\vfill % Fill the rest of the page with whitespace

\end{titlepage}


\begin{interlude}
\listoftodos
\end{interlude}


\cleardoublepage% especially in a document where chapters start at right-hand pages
%\phantomsection% for an anchor if you use hyperref

\chapter*{Remerciements}% for the actuall unnumbered heading
\addcontentsline{toc}{chapter}{\numberline{}Remerciements}% if you wish to have a TOC entry
\thispagestyle{empty}% or plain etc.
\markboth{Remerciements}{Remerciements}% relevant depending on page style
% or if it's more than one page


Je tiens à remercier les personnes qui m'ont permis de prés ou de loin à accomplir ce stage et qui m'ont aidé lors de la rédaction de ce rapport.

Tout d'abord, j'adresse mes remerciements à Messieurs \textbf{Fabrice Mahé et Eric Darrigrand} profeseurs de l'Université de Rennes 1, qui m'ont permis de suivre cette formation et qui m'ont accompagné lors de la recherche de stage.

Je tiens à remercier Monsieur \textbf{Geoffroy ETAIX}, qui m'a accordé sa confiance pour le stage et pour etc 
\todoCompleter{remercier Goeffroy}
\todoRediger{Remercier Fongecif}
\todoRediger{Remercier Ingeus}
\todoRediger{Remercier IDApps, Yann, Stéphane, et surtout Laila!}

Enfin, last but not least, je tiens à remercier mon épouse Céline et mes deux filles pour leur encouragments, leurs soutient et leur patience.
\todoCompleter{Céline et les filles}

\todoCompleter{Completer les remerciements}
\todoMiseEnForme{Aérer}


\setcounter{tocdepth}{5}
\tableofcontents
\newpage

\chapter{Introduction}
\todoAmeliorer{L'introduction doit-elle être un chapitre ou alors un Abstract?}
\todoRediger{Introduction du context projet: j'ai effectué le stage blabla au sein de l'entreprise bla etc}
	\section{Tellus Environment - Missions Principales}
	\todoRediger{Tellus création, mission, activité}
	
	\subsection{Géophysique et cartographie haute-définition}
	Tellus Environmenent est une startup spécialisée dans la cartographie haute définition des sous-sols et des fond-marins. Elle propose à ses clients une offre bout en bout d'acquisition et de traitement des données pour permettre à ses clients d'agir en fonction de ces conclusions.
	\todoAmeliorer{Ce \S $\;$peut servir à l'intro sur Tellus.}
	
	\subsection{Collecte des données géophysique}
	
	\paragraph*{} Tellus Environment met en oeuvre ses propres équipements géoradar et magnetomètre pour l'acquisition de données sur le sous-sol sur des surfaces de quelques mètres-carré à quelque hectares. Elle est aussi en mesure de planifier et coordonner la mise en oeuvre d'équipements plus lourds - LIDAR, géoradar, magnétomètres aéroporté - pour obtenir des données sur des surfaces beaucoup plus importantes, de l'ordre du kilomètre carré.
	
	\missingfigure{Georadar}
	\missingfigure{Magnetomètre}
	\missingfigure{Radargramme}
	
	\paragraph*{} Tellus Environment peut aussi coordonner la mise en oeuvre d'équipements d'acquisition marins - sonar, multibeam, \underline{percuteur} \todoCheck{est-ce le bon terme?} - pour permettre la cartographie des fonds marins et des sous-sols aquatiques, que ce soit en environnement eau douce - rivières, lacs étangs - ou marins. 
	\missingfigure{sonar}
	\missingfigure{sonargramme}
	
	\paragraph*{} Tellus Environment a accès à de nombreuses bases cartographiques pour compléter les données acquises sous sa supervision pour aider à la mise en oeuvre des équipements ainsi qu'à compléter les données acquises en vue de leur traitement.

	\subsection{Traitement des données}
	\todoRediger{Expertise Tellus Traitement des données}
%	\todoObjectif{Montrer la pertinance de l'expertise de Tellus pour la cartographie 3D}

	\section{Activité R\&D}
	\todoRediger{Activité R\&D}
	
	Activité: R\&D: développer des composants qui mettent en oeuvre les expertises de TellusEnvironment pour créer des produits innovants, dont les composants seraient de plus réutilisables pour améliorer la productivité des services Géophysique.


	\section{Projet Symeter V2}
	\todoRediger{Context Projet}
		\subsection{Symeter V1}
		Symeter V1: rapide rappel
		\todoRediger{SymeterV1}
		
		\subsection{Symeter V2: Objectifs}
		\todoRediger{Symeter V2: objectifs}
		Symeter V2: Objectifs
		Symeter V2: Présentation du plan de projet
- test du simulateur Gazebo pour évaluer son utilité dans le projet Symeter2
- Simulation couverture lidar
- Montage des outils nécessaires au développement simulé du projet symeter
- Mise en place de la localisation: installation et tests
- Mise en place de l'acquisition des relevés.

\chapter{Technologies, Contraintes et Plan de Projet}
	Ce chapitre présente dans leurs grandes lignes les entrants et les contraintes mis en jeu par le projet Symeter V2. On y énumère notament les différents éléments à mettre en oeuvre pour monter le système Symeter V2: les capteurs l'environnement d'exploitation logicielle, les outils mathématiques et procédés à coordonner pour monter le prototype de Symeter 2.
	
	\section{Capteurs}
		Les capteurs à mettre en oeuvre seront un ou plusieurs LIDAR, un ou plusieurs IMU et un ou plusieurs GPS. Cette section décrit les fonctionnalités de ces capteurs de ces capteurs.
			
		\subsection{LIDAR}
		Un LIDAR est un équipement qui permet de prendre de nombreuses mesures de l'environnement à partir d'un laser. En exploitant les informations d'angle de pointage et en observant les caractéristiques du signal réfléchin tant du point de vu temporel (temps de vol, phase) que du point de vue de la forme du signal (puissance réfléchie, déphasage, distorsion, effet Doppler , ...), il est possible de mesurer de nombreuses caractéristiques de l'environnement.
		
		\paragraph*{} Le LIDAR est utilisé par Symeter V2 pour détecter la forme précise du chantier d'ensilage.
		
		\subsubsection{Détecteur Actif:}
		Le LIDAR est essentiellement un détecteur actif dont de le principe est identique à celui du radar: un pulse electromagnétique est émis dans une direction privilégiée par une partie émettrice, constituant un signal incident. 
		
		\para Ce signal incident est réfléchi par un object de l'environnement, générant un signal réfléchi qui retourne en direction du LIDAR. Ce signal réfléchi est reçu par un détecteur. 
		
		\para En comparant le signal émis et le signal reçu nous pouvons en déduire certaine caractéristiques de l'environnement.
		
		\paragraph*{Différents types de mesure: } Il existe deux grands types de LIDAR: \begin{description}
			\item[LIDAR  à "temps de vol":] les lidars de ce type déterminent la distance de la cible en mesurant le temps écoulé par entre le moment de l'émission du pulse et la réception du signal réfléchi par la cible.
			\item[LIDAR  "full waveform":] les LIDAR de ce type émettent un signal périodique continu et mesurent la distance de l'object détecté essentiellement par mesure de déphasage du signal réfléchi par rapport au signal émis. Les modifications de la forme du signal réfléchi et de son intensité permettent d'extraire d'autre informations à propos de la cible. 
		\end{description} 
		
		\subsubsection{Scan de l'environnement}
		En faisant varier la direction du bloc émetteur/récepteur du LIDAR, il est possible de faire plusieurs mesures de l'environnement. Le LIDAR mettant en oeuvre des procédés optiques, il est possbile de faire varier très rapidement son pointage et ainsi d'obtenir de très nombreux relevés de l'environnement.
		
		\paragraph*{Deux grands modes de scan de l'environnement:} Il existe deux grands modes de scan de l'environnement pour un LIDAR
		
		\begin{description}
			\item[Mode 2D ou  "planar":] Dans ce mode le LIDAR mesure l'environnement dans un plan unique, généralement en faisant tourner le faiseau de manière circulaire autour d'un axe.
			\item[Mode 3D:] Dans ce mode le LIDAR mesure des portions d'espace, généralement en combinant deux rotations simulatées du faiseau, soit en faisant tourner sur un même axe de révolution plusieurs faiseaux.
		\end{description}
	
	\missingfigure{mode 2D vs mode 3D}
		
		\subsubsection{LIDAR Hokuyo UTM-30LX-EW}
		Le projet Symeter V2 met en oeuvre un LIDAR de type Hokuyo UTM-30-LX-EW, le même que celui utilisépar Symeter V1. C'est un équipement robuste et précis qui est de plus bien connu par Tellus Environment.
		
		\para Cet équipement est basé sur un laser infrarouge (longueur d'onde de 905nm) pour scanner un champ semi-circulaire de 270$\degres$. C'est un LIDAR de type "planar, temps de vol".
		
		\para Il mesure la distances des objets à sa portée par pas de 0,25$\degres$. La distance de détection maximale est de 30m.
		\newline
		
		\begin{tabular}{|c|c|}
			\hline 
			Angle de Balayage & 270$\degres$ \\ 
			\hline 
			Résolution Angulaire & env. 0,25$\degres$ \\ 
			\hline
			Temps de Balayage & 25ms/balayage (40 balayages par seconde) \\
			\hline
			Distance de détection & Portée garantie: 0,1 à 30m \\
			\hline
			Résolution de la mesure & 1mm \\
			& $\sigma$ \\
			\hline
		\end{tabular} 
	
		\todoRediger{Décrire les caratéristiques exactes }
		\missingfigure{Photo du Hukyo}
		
		\subsection{IMU}
		Un IMU - Inertial Mesurement Unit - est un equipement comportant plusieurs capteurs inertiels, d'accélération et d'angles qui permettent avec un traitement des données approprié d'établir de suivre la pose du véhicule sur lequel il est monté.
		
		\subsubsection{Les grandeurs mesurées}
		Un IMU comporte généralement 3 accéléromètres linéaires et un gyroscope, et souvent un magnétometre.
		
		\para Le \textbf{gyroscope} permet de mesurer, avec précision et en continu, l'inclinaison du véhicule en terme de \textbf{roulis, tanguage et lacet}. 
		
		\para Les \textbf{accéléromètres} mesurent en continu les accélérations linéaires de la partie du véhicule sur lequel l'IMU est fixé, dans trois directions différentes. Ces directions sont généralement dénotées $x,y,z$, par analogie avec le repère cartésien. Simplement, les accélérations sont mesurées dans le repère du véhicule, mobile et variable et non dans un repère absolu immobile.
		
		\para En combinant les variations d'inclinaison et les mesures d'accélérations linéaires dans le repère mobile du véhicule, il est possible de reconstruire avec une certaine certitude la trajectoire du véhicule dans le repère absolu, en utilisant les équations de dynamique Newtonnienne.
		
		\missingfigure{Trajectoire et mesures}
		
		\para Comme indiqué ci-dessus, un IMU possède souvent un \textbf{magnétomètre}. Celui-ci permet de mesurer la direction et l'intensité du champ magnétique terrestre par rapport à l'IMU et permet d'en dériver le cap, la direction vers lequel le véhicule pointe sur la surface de la Terre. 
		
		\para Ce permet, en démarrant l'IMU alors que le véhicule est parfaitement immobile, d'obtenir les conditions initiales du vecteur d'état du véhicule, par rapport au repère odométrique.
		
		\para Par ailleur, reconstituer une position à partir d'informations d'accélération implique mathématiquement de procéder à une double intégration numérique. Ce procédé est donc sujet à une dérive dans le temps et doit donc être recalé par une mesure de pose issue d'un autre capteur.
		
		\para Par contre, un IMU peut fournir ses informations avec un très haut taux de rafraichissement, sans nécessiter d'information extérieure au véhicule. Il peut donc se révéler indispensable lorsqu'aucune autre nouvelle information sur la position du véhicule n'est disponible sur des périodes de temps plus ou moins longue.
		
		\subsubsection{IMU XSense}
		\todoRediger{données générées par un IMU}
		\missingfigure{graphes temporels des données générées par l'IMU}
		\missingfigure{Photo de l'IMU xsense}
		
		\subsection{GPS en mode RTK}
		
		Le système Symeter V2 ayant besoin d'une précision centimétrique pour permettre la mesure du tas d'ensilage avec précision du même ordre de grandeur, il mettra en oeuvre d'un GPS en mode RTK. Ce mode fonctionnement particulier du système GPS est en mesure de fournir position centimétrique du véhicule relativement à une borne GPS (base) situé à proximité du chantier.
		
		\subsubsection{GPS simple}
		Pour expliquer le principe du GPS en mode RTK, il est nécessaire de comprendre comment le système GPS "simple" fonctionne. 
		
		\para Le système GPS est un système de positionnement par satellite qui, en se basant sur les signaux émis par une constellation de satellites, est en mesure de calculer la position sur la surface de la Terre d'un récepteur avec une précision de l'ordre de la dizaine de mètres, pour une utilisation courante.
		
		\para Le principe de base est celui de la triangulation: les satellites GPS émettent en permanence un signal comportant une horloge (temps GPS). En connaissant les éphémérides des satellites et en écoutant les signaux de plusieurs de ces satellites, notamment leur signal d'horloge respectif, un récepteur GPS peut calculer avec une grande précision la position de chacun des satellites émetteurs dans le référentiel géocentrique.
		
		\para Une fois le positionnement des satellites établi, une simple triangulation permet de calculer avec une grande précision du récepteur sur la surface de la Terre.
		
		\para Pour obtenir une précision suffisante, un récepteur doit écouter au moins 4 satellites en même temps.
		
		\missingfigure{GPS en mode simple}
		
		\subsubsection{Principe du mode RTK}
		Le mode RTK du système GPS est un mode de fonctionnement qui permet de mesurer avec une très grande précision la position relative de deux récepteurs GPS distincts. L'un, dénommé \textbf{la base} est un récepteur fixe, dont la position GPS est connue avec une grande précision. L'autre dénommé \textbf{le rover} est en communication constante avec la base au moyen d'un liaison radio adaptée.
		
		\para Les deux récepteurs étant présents sur le même site à la surface de la Terre, ils seront à l'écoute tous les deux mêmes satellites. En comparant les signaux reçus d'un même satellite par la base et le rover, le rover est en mesure de calculer avec une très grande précision - de l'ordre du centimetre - la différence de distante entre la base et le satellite d'une part, et le rover et le satellite de l'autre part. 
		
		\missingfigure{GPS en mode RTK}
		
		\para En combinant ces différences de distance à partir de plusieurs satellites, le rover est donc capable de calculer très précisément sa position relativement à la base.
	
		\todoTrouverRefExt{GPS en mode RTK}
		
		
		\missingfigure{Photo base + rover RTK}
		
		\todoPeutEtre{caméra video pour odométrie visuelle}
		
	\section{Environnement de programmation: ROS}
	Du fait de son utilisation dans la première version de Symeter, et du fait de ses qualités en terme de modularité, de disponibilité des drivers pour les capteurs et actionneurs, l'environnment logiciel ROS a été choisi avant même le début du stage pour être la base de la partie logicielle de Symeter V2.
	
	\missingfigure{Quelques robots powered by ROS}
	
	\para ROS est un environnement logiciel destiné à la mise en oeuvre de plateforme robotique. Cet environnement très modulaire permet de programmer et de déployer de nombreux modules interdépendants, appelés \textbf{packages}, constitués eux de noeuds exécutables appelés \textbf{nodes}.
	
	\para Ces nodes peuvent communiquer entre eux aux moyens de deux mécanismes principaux: \begin{itemize}
		\item les services
		\item les topics
	\end{itemize}

	\para Ces deux mécanismes permettent l'échange de données entre nodes selon des messages de la structure peut être spécifiée par le développeur du système.
	
	\para Les \textbf{services} sont mis à disposition chacun par un node. Ils peuvent être invoqués par d'autre node en effectuant une transaction explicite: le node appelant fourni une donnée d'entrée au node mettant à disposition le service. Ce dernier traite la donnée entrante puis génére une donnée de sortie qui est enfin renvoyée au node, terminant ainsi la transaction.
	
	\para Les \textbf{topics} sont des canaux de diffusion d'information similaire à une autoroute: tout node peut diffuser un message vers un topic et tout node peut écouter un topic pour recevoir et traiter les messages qui y sont émis.
	
	\missingfigure{Exemple d'une structure en packages, noden topics et services}
	
	\para Ces topics pourvoient des données que l'on pourrait qualifier de "sensorielle", dont la durée de vie est courte, telle que les données générées par un capteur fournissant des données en continu. 
	
	\para Ils sont aussi le support de diffusion des \textbf{transformations}, qui permettent de décrire l'état de \textbf{pose} de chacuns des composants mécaniques d'un véhicule autonome, ainsi que la pose du véhicule dans son environnement.
	
	\para ROS est un environnement fourni sous une licence Open Source et propose des drivers pour de très nombreux capteurs et actionneurs de tous types. Il intègre aussi des composants logiciels tiers spécialisés tels que, entre autre, la Point Cloud Library pour le traitement des nuages de points et la librairie OpenCV pour le traitement de la vision par ordinateur.
	
		
	\section{Les outils mathématiques}
	\todoObjectif{Introduire les concepts, indiquer qu'ils seront approfondis dans la suite du document.}
	
	\para Comme indiqué en introduction, le système Symeter V2 sera une plateforme semi-autonome, qui devra être en mesure de fournir une aide à la décision à partir de mesures effectuées depuis une plateforme mobile, et ceux avec une intervention humaine minimale.
	
	\para Symeter V2 devra donc estimer en temps réel sa pose, afin de pouvoir intégrer une représentation 3D du chantier afin de pouvoir effectuer son service.
	
	\para Les outils mathématiques qui devront être mis en oeuvres sont essentiellement ceux mis en oeuvre dans les systèmes robotiques:
	 \begin{itemize}
	 	\item Les \textbf{poses}, c'est à dire la représentation de la position des capteurs et des éléments mécaniques du véhicule les uns par rapport aux autres, et celle du véhicule par rapport à son environnement.
	 	\item L'estimation de paramètres à partir d'un ensemble de mesures issues de différents capteurs, au moyen de la \textbf{fusion de données}
	 	\item Reconstruction d'une scène 3D à partir de nuages de points
	 \end{itemize} 
 
   \para Ces outils sont décrits plus en détail dans les section suivants.
	
	
		\subsection{Positionnement en Robotique: Poses}
		
		La robotique à pour enjeu de permettre la mise en oeuvre de systèmes mécaniques autonomes, mobiles ou non, qui mettent en oeuvre une série de capteurs et d'actionneurs pour agir sur leur environnement avec une interaction humaine très limitées.
		
		\para La capacité de représenter de manière fine la position des différents éléments d'un tel système les uns par rapport aux autres, que ce soit en position $x,y,z$, mais aussi en inclinaison - $\theta_x$ (tanguage), $\theta_y$ (roulis), $\theta_z$ (lacet) - est donc primordiale. L'état de l'élément décrit par ce vecteur $x,y,z, \theta_x, \theta_y, \theta_z$ est appelé sa \textbf{pose}.
		
		\missingfigure{Pose 3D d'un objet}
				
		\para Les outils mathématiques de prédilection pour décrire une pose sont bien sur les transformations dans $\mathbb{R}^3$, translations et rotations. Cependant le système de coordonnées classique à 3 dimensions spatiales $x,y,z$ ne permet pas de combiner simplement les transformation 3D, ce qui est un besoin de base pour passer d'un réferentiel à un autre.
		
		\para C'est pour cela que la robotique fait une utilisation intensive des \textbf{Coordonnées Homogènes}, issues de la géométrie projective, qui permettent de combiner les transformation 3D en passant par un espace à 4 dimensions approprié. Les transformations y sont représentées par des matrices de $\mathbb{R}^{4 \times 4}$, ce qui permet de combiner les transformations par simple multiplication de matrices.
		
		\para Les principes de base des coordonnées homogènes sont décrites en Annexe XXX \todoAjouterRef{Lien vers l'annexe}
		
		\para Les rotations dans l'espace sont aussi représentables par des quaternions de norme unitée. Cette représentation est souvent préférées par la robotique par rapport à une représentation basée sur les rotations eulériennes car elle permet de combiner les rotations sans rentrer dans un "gimbal lock". 
		
		\para ROS utilise les quaternions de manière native dans la description de ses transformations. \todoAjouterRef{Quaternion et quaternions unité}
		
			
		\subsection{Localisation par fusion de données}
		
		Afin de pouvoir reconstituer le chantier en une représentation 3D avec une précision centimétrique, il est nécessaire que pour chaque mesure LIDAR le système Symeter V2 aie une estimation suffisamment précise de la pose du véhicule, afin de déterminer la pose du LIDAR et pouvoir situer dans le chantier la position $x,y,z$ de chaque point mesuré.
		
		\para Le problème de la localisation du véhicule est donc le suivant:
		\para Estimer à chaque instant $t_n$ le vecteur d'état:
		\begin{equation}
		\mathbf{X} = \begin{bmatrix}
			x_n \\ y \\ z \\ \theta_x \\ \theta_y \\ \theta_z
		\end{bmatrix}
		\end{equation}
		du véhicule à partir des données disponibles.
		
		\missingfigure{Flux de données IMU + GPS --> $X$}
		
		\para Les sources de données disponibles pour estimer la pose du véhicule à chaque instant sont:
		\begin{itemize}
			\item L'IMU,
			\item Le GPS en mode RTK.
		\end{itemize}
		
		\para L'IMU fourni des informations sur l'inclinaison du véhicule, ainsi que les accélérations linéaires qu'il subit à une fréquence de l'ordre de 30 Hz. Cependant comme indiqué ci-dessus, calculer la position du véhicule à partir de la mesure des accélérations subies requiert une double intégration, sur un signal d'accélération qui estgénéralement très biruité. Ces éléments impliquent une forte probabilité de dérive dans le temps.
		
		\para Le GPS en mode RTK peut fournir une mesure de la position du véhicule dans la scène avec une précision de l'ordre de quelques centimètres, mais à un taux de rafraichissement beaucoup plus bas, de l'ordre de 2 à 5 Hz. De plus le GPS ne fournit aucune information sur l'assiette du véhicule.
		
			
		\todoRediger{introduire les filtres de Kalman}
		\missingfigure{chaine de traitement de fusion de données}
		\todoAjouterRef{Filter de Kalman pour la fusion de données.}
		\todoCompleter{Comme le GPS fourni une position absolue dans le référentiel de la scène nous ne sommes pas dans une problématique SLAM. Le véhicule n'est pas autonome, la navigation est assurée par l'opérateur du tracteur.}
			
		
		\subsection{Traitement des nuages de points}
		Le LIDAR utilisé par le système Symeter V2 étant de type "2D planar", les points retournés lors d'une unique mesure feront tous partie du même plan. L'ensemble des mesures sera donc constitué d'un ensemble de tranches qu'il faudra reconstituer à l'aide d'un traitement adapté.
		
		\para Le traitement des données LIDAR comportera les étapes suivantes:
		\begin{enumerate}
			\item Filtrage des données pour une obtenir une densité de points de mesure qui soient suffisament espacés
			\item Convertion des données LIDAR (direction / distance) en un nuage de points (liste de point $P_n$ avec chacun des coordonnées $x_n, y_n, z_n$
			\item Stockage de ces points dans une structure de données capable d'accumuler les nouveaux points mesurés au cours du temps, et constituer une carte 3D d'occupation.
 		\end{enumerate}
 	
 		\para La plupart de ces traitements seront implémentés en utilisant le Point Cloud Library qui offre de nombreux composants logiciels intégrés directement dans ROS.

		\missingfigure{Architecture de traitement des nuages de points}
			
	\section{Contraintes de développement}
		Le système Symeter V2 st destiné à évoluer sur un véhicule de type tracteur agricole. Il est de plus destiné à construire une représentation 3D d'un chantier d'ensilage, partir de positions qui vont varier non seulement en $x$ et en $y$, mais aussi et surtout en $z$.
		
		\subsection{Capacités de tests en grandeur limitées}
		Pour développer et tester le système Symeter V2, le développeur ne disposera pas de tracteur en grandeur réelle, ni de chantier d'ensilage, ceux-ci étant des chantiers annuels se déroulant à des moments bien défini dans l'année (printemps et automne).
		
		
		\para Par ailleurs, au début du stage il n'est pas encore tout à fait décidé de combien de LIDAR seront nécessaires pour que le système puisse assurer une couverture complète depuis des équipements embarqués sur le tracteur.
		
		\para Les équipements à mettre en oeuvre sont relativement couteux et leur mise en oeuvre requiert une certaine expertise. Il est donc nécessaire de pouvoir mettre en oeuvre le développement du système par le biais de mises en oeuvre alternatives.
		
		\para Ces possibilités sont au nombre de 2: la simulation et la mise en oeuvre en grandeur des équipements sur une camionnette pour effectuer des tests simples de reconstitution du terrain.
		
		\subsection{Plateformes de test disponibles}
			
			\subsubsection{Environnement de simulation Gazebo}
			
			Gazebo est un environnement de simulation très complet, qui permet de simuler un environnement physique dans lequel une plateforme semi-robotique simulée peut évoluer. 
		
			\todoRediger{De nombreux capteurs virtuels disponibles, possibilité de faire varier les incertitudes}
			
			\para Gazebo est un logiciel intégré à ROS. Sa mise en oeuvre requiert de programmer un véhicule virtuel similaire à la plateforme cible du système (dans notre cas un tracteur avec les capteurs listés dans les sections précédentes), puis de déployer le logiciel ROS pour Symeter V2 sur ce véhicule simulé. 
			
			\para Il est aussi possible de spécifier le "monde" dans lequel le tracteur simulé évoluera, en y ajoutant des murs, du relief, voire une butte en guise de tas d'ensilage.
		
			\missingfigure{Capture d'écran Gazebo avec le tracteur et le silo}
			
			\para Gazebo est donc utilisé pour simuler un chantier d'ensilage pour tester la localisation, l'acquisition du terrain, la mesure d'un modèle de Silo.
			
			\subsubsection{Prototype monté sur Camionnette}
			
			Une simulation n'étant jamais parfaite (comme nous le verrons par la suite), il est quand même nécessaire d'effectuer des certains tests avec des instruments réels, lors de tests en grandeur.
		
			\missingfigure{Maquette avec équipements réels montés sur une camionnette.}
	
				
			\para Pour la mise en oeuvre en grandeur, un prototype de l'équipement a été monté sur la camionnette de Tellus Environnement, la "Tellus Car".
			
			\para Le prototype comporte:
			\begin{itemize}
				\item Une unité de calcul
				\item Un GPS en mode RTK
				\item Un LIDAR de type Hkuyo UTM-30LX-EW
				\item Un IMU de type XSense
			\end{itemize}
		
			\para Ce prototype a été utilisé pour effectuer des captures des trois capteur alors que la Tellus Car évoluait dans le parking adjacent aux locaux de Tellus Environnement.
		
	\section{Les grandes parties du projet}
	
	\todoRediger{comme il est apparent, proejt avec de nombreux sujets, et nous partons de zéro.}
	\todoRediger{objectif est donc de monter une plateforme de base avec dans un premier temps et autant que possible des composants "sur étagère", d'en évaluer les performances et de dégager les axes d'améliorations nécessaires.}

	
	\todoRediger{Montage de la plateforme simulée ROS + Gazebo}
	\todoRediger{Elaboration, prototypage et test du processus de localisation}
	\todoRediger{Elaboration, prototypage et test du procédé de capture de la scène 3D}
	\todoRediger{Test en grandeur à l'aide d'un équipement monté sur la Tellus Car.}
				

\chapter{Simulation d'un tracteur évoluant sur un chantier d'ensilage à l'aide de ROS/Gazebo}

\todoRediger{introduction du chapitre simulation tracteur}

	\section{Présentation de ROS}
	\todoRediger{Présenter succinctement ROS, montrer quelques projets, exposer les outils}
	\missingfigure{Photo d'exemples de robots, roues, bras, humanoides, etc}
		\subsection{Architecture de ROS}
		\todoRediger{ROS est un environnement permettant le montage de plateformes robotiques complètes.}
		\todoRediger{Architecture logicielle très modulaire}
		\todoRediger{Nodes, services, topics, capteurs, etc}
		\missingfigure{Exemple de node}
		\subsection{Gestion des transformations}
		\todoRediger{précision d'évolution d'un plateforme robotique dépend de la prise en compte de la position relative de ses différents capteurs et actionneurs}
		\subsection{Gestions des Capteurs}
		
		\todoRediger{un petit laïus surt les capteurs supportés}
	\section{Présentation de Gazebo}
	
	\todoRediger{Qu'est ce que Gazebo, à quoi ça sert}
		\subsection{Construction d'un robot virtuel}
		\todoRediger{Décrire les éléments nécessaires à la description du robot}
		\todoRediger{URDF et SDF, world, contrôleurs}
		
		
		\subsection{Vérification de disponibilité des capteurs}
		\todoRediger{Besoin de vérifier la disponibilité et la performance des capteurs cibles de Symeter V2}
		
	\section{Contraintes de mise en oeuvre}
		\subsection{Quelques bugs génants}
		\todoRediger{Lister les problèmes et bugs qui ont nécessité la mise en place de contournements}
		\subsection{Simulation mécanique, frottements, adhérence}
		\todoRediger{La nature des "matériaux" utilisés dans gazebo est très "rigide": les chocs ne sont pas amortis.}
		\subsection{Conclusions sur les contraintes}
		Eviter à tout prix que le véhicule entre en dérapage: il ne récupère généralement pas.
		Raisons pour lesquelles le véhicule peut entrer en dérapage:
			- mauvaise configuration des paramètres de frottement des éléments en contact avec le sol.
			- vitesses de rotation différentes des roues: 
				- mauvaise consigne de vitesse sur les roues
				- différence de vitesse entre roue intérieure et roue extérieure lors d'un virage
			- en virage: rayons de virages différent entre roue intérieure et roue extérieure
			
		Conduite en terrain accidenté:
			- Conserver la directivité
			- 4 roues motrices et/ou suspension
			
			\todoCompleter{Mise en forme}
		
	\section{Mise en Oeuvre: simulation d'un environnement de tassage de silo}
		\subsection{Montage d'un tracteur simulé}
		Intérêt: simuler l'implantation physique des capteurs avec des dimensions du même ordre de grandeur que les plateformes cibles.
		
		\subsubsection{Description Physique}
			\paragraph{Chassis}
			Le chassis du tracteur à été monté en se basant sur les dimensions générales d'un tracteur de type \textbf{TODO: ICI TYPE DE TRACTEUR} \todo{Ecrire type de tracteur}, en utilisant cependant des formes géométriques simplifiées.
			\newline
			
			Le tracteur est donc composé d'un pavé en guise de chassis, de 4 cylindres allongés en guise d'essieux, et de 4 grands cylindres en guise de roues.
			
			 \missingfigure{Image du tracteur modélisé}
			\paragraph{Actuateurs et Contrôleurs}
			\todo{Rédiger Actuateur et Contrôleurs}
		\subsubsection{Propulsion et Guidage}
			\paragraph{Algorithme}
			
			\todo[inline]{Description géométrique du problème de direction différentielle}
			
			\paragraph{Implémentation sous ROS}
			\todo{Description de l'implémentation sous ROS}
			\missingfigure{Schéma des flux topic qui permet de transformer une commande twist en 4 commandes de velocités.}
			
		
		

\chapter{Mise en place du processus de localisation}

	\section{Présentation du problème}
	\todoRediger{Exposer la problématique: fusionner des données avec potentiellement des fréquences d'échantillonnage différentes}
	
	\todoRediger{déterminer la pose du tracteur à partir des différents capteurs}
	
	\section{Filtres de Kalman}
	\todoRediger{théorie des filtres de Kalman}
	
	\subsection{Filtres de Kalman Linéaires}
	
	\subsection{Filtres de Kalman Etendus}
	\todoRediger{pour les systèmes dynamiques non linéaires}
	
	\section{Mise en oeuvre}
	
	\todoRediger{Utilisation du module robot\_localization}
	
	
\chapter{Exploitation des données LIDAR}

	\todoRediger{Exposer le plan du chapitre}
	\todoRediger{quelle est la qualité de mesure d'un LIDAR sur une surface potentiellement inclinée?}

	\section{Présentation de la chaine de traitement des données LIDAR}
	
	\missingfigure{Chaine de traitement LIDAR --> Nuage de Point}
	\missingfigure{Transformer un ensemble de lignes en un nuage de points cohérent}
	
	\missingfigure{Inclinaison LIDAR par rapport au tracteur}
	
	\todoRediger{Biblio sur l'influence de l'inclinaison du lidar par rapport à la surface}
	\todoRediger{arg1}
	
	\section{Acquisition des données LIDAR}
	

	\todoRediger{mettre en forme les données lidar de manière à pouvoir les utiliser}
	\todoRediger{réduire la bande passante, ne garder que les données exploitatables}
	
		\subsection{Transformation trame LIDAR en un nuage de points}
		\todoRediger{LIDAR: ensemble d'angles, temps de vol, puissance reçue}
		\todoRediger{Nuage de points: positionnement x,y,z dans la scène}
		\todoRediger{Il est nécessaire de connaitre précisement la pose du LIDAR pour générer le nuage de points.}
		
		\missingfigure{Illustration nature donnée LIDAR}
		\missingfigure{Illustration nature donnée Nuage de Point}
		\missingfigure{Diagramme fonctionnel LIDAR --> Nuage de Point}
		
		\subsection{Filtrage de la ligne par traitement Voxel}
		\todoRediger{Enjeu: réduite la bande passante, ne pas sur-échantillonner au pied}
		
		\missingfigure{suréchantillonage au pied du tracteur, sous-échantillonnage sur les côtés.}
		
		\todoRediger{Exposer le principe du filtrage}
		
		\todoRediger{mise en oeuvre grace à la point cloud library}
		
	\section{Accumulation des nuages de points}
	\todoRediger{Objectif: reconstituer la scène à partir des nuages de points "ligne" obtenus à partir du processus d'acquisition}
	
		\subsection{Le principe de fonctionnement}
		\todoRediger{Garder en mémoire les nuages de points correspondant à des lignes précédemment acquise}
		\todoRediger{ne pas dupliquer les informations: éliminer les points doubles}
		
		\subsection{Principe de stockage des données 3D}
		\todoRediger{les structures de données de stockage de l'information 3D}
		\todoTrouverRefExt{les structures de données 3D}
		
		\subsubsection{Les B-Trees}
		\todoRediger{Les b-trees, principes}
		\todoRediger{Point forts: calculs performants}
		\todoRediger{points faible: arbre équilibrés --> difficile d'ajouter de nouveaux points.}
		
		\subsubsection{Les Octrees}
		\todoRediger{Octree, principes}
		\todoRediger{points forts: structure déséquilibrées sans problème, possibilité d'ajouter de nouveau points avec une bonne performance}
		\todoRediger{point faible: beaucoup d'overhead de memoire si pas implémenté correctement.}
		
		\subsubsection{Le choix: octree correspond à notre besoin.}
		
		\subsection{Mise en oeuvre: octomap}
		\todoRediger{Utilisation d'octomap car composant sur étagère}
		\todoRediger{permet d'implémenter rapidement la chaine de traitement pour vérifier la validité de la faisaibilité}
		
	\section{Mise en oeuvre sous Gazebo}
	\todoRediger{implémentation d'un monde avec un silo et un tas d'ensilage}
	\todoRediger{capture à l'aide d'un lidar monté à l'arrière du tracteur}
	
		\subsection{Test sous gazebo}
		\todoRediger{Tracteur effectue une acquisition initiale de la scène}
		\todoRediger{Une fois l'acquisition effectuer, on invoque le node de sauvegarde de nuage, qui le sauve dans un fichier}
		
		\subsection{Analyse du nuage de point généré}
		\subsubsection{Outil pour l'analyse de nuage de points: Paraview}
		\todoRediger{Outil de visualisation utilisé: Paraview}
		\todoRediger{besoin de convertir pcd en vtk à l'aide de l'outil approprié}
		\todoRediger{permet de naviguer de manière efficasse dans le nuage de point, d'effectuer des projections, etc, etc}
		
		\subsubsection{Points à améliorer sur le nuage de points}
		\todoRediger{Le nuage est épais}
		
		\todoRediger{octomap génère un nuage de point basé sur le centre des voxel occupés --> problème de quantification volumique implique perte de précision de la mesure}
		
		\todoRediger{Soit réduire la taille du voxel --> augmentation de la taille}
		
		\todoRediger{Soit retourne le point moyen de chaque voxel --> non supporté par octomap --> à implémenter nous même.}
		
		\todoAjouterRef{octomap}
		

\chapter{Mise en oeuvre à partir de mesures réelles}

	\section{Protocol de test}
	
	\section{Données générée - visualisation sous google maps}
	
	\section{Exploitation des données}
		\todoRediger{données sans RTK --> besoin d'avoir un}

\chapter{Reste à faire et Améliorations}


\begin{appendix}
	\chapter{Positionnement En Robotique}
		\section{Géometrie projective, Coordonnées Homogènes}
		\section{Une autre descriptions des rotations en 3D: Quaternions Unitaires}
		\section{Application: Simulation de couverture d'un faiseau LIDAR orienté vers le sol}
	
	\chapter{Filtres de Kalman}
	
	\chapter{ROS: Architecture et Concepts}
	
	\chapter{Point Cloud Library}
	Essai de citation \cite{kaplan_understanding_2006}.

\end{appendix}

\nocite{*}
\bibliographystyle{alpha}
\bibliography{biblio}

\end{document}